% Part: first-order-logic
% Chapter: axiomatic-deduction
% Section: proof-theoretic-notions

\documentclass[../../../include/open-logic-section]{subfiles}

\begin{document}

\iftag{FOL}
      {\olfileid[fr]{fol}{axd}{ptn}}
      {\olfileid[fr]{pl}{axd}{ptn}}

\olsection{Notions de théorie de la démonstration}

\begin{explain}
De même que nous avons défini plusieurs notions sémantiques importantes
(\iftag{FOL}{validité}{tautologie}, conséquence logique et
satisfaisabilité), nous définissons maintenant les \emph{notions de
théorie de la démonstration} correspondantes. Celles-ci ne font pas
appel à la satisfaction d'!!{sentence}s dans des !!{structure}s, mais
à la !!{derivability} ou à la non-!!{derivability} de certaines
!!{formula}s. La coïncidence de ces notions constitue une découverte
importante : c'est le contenu des \emph{théorèmes de correction et de
complétude}.
\end{explain}

\begin{defn}[!!^{derivability}]
!!^a{formula}~$!A$ est \emph{!!{derivable} à partir de~$\Gamma$}, ce
que l'on note $\Gamma \Proves !A$, s'il existe une !!{derivation} à
partir de~$\Gamma$ qui se termine par~$!A$.
\end{defn}

\begin{defn}[Théorèmes]
!!^a{formula}~$!A$ est un \emph{théorème} s'il existe une
!!{derivation} de~$!A$ à partir de l'ensemble vide. On écrit
$\Proves !A$ si $!A$ est un théorème, et $\Proves/ !A$ dans le cas
contraire.
\end{defn}

\begin{defn}[Cohérence]
Un ensemble~$\Gamma$ de !!{formula}s est \emph{cohérent} si et seulement
si $\Gamma \Proves/ \lfalse$; il est \emph{incohérent} dans le cas
contraire.
\end{defn}

\begin{prop}[Réflexivité]
\ollabel{prop:reflexivity}
Si $!A \in \Gamma$, alors $\Gamma \Proves !A$.
\end{prop}

\begin{proof}
La !!{formula}~$!A$ constitue à elle seule une !!{derivation} de~$!A$
à partir de~$\Gamma$.
\end{proof}

\begin{prop}[Monotonie]
\ollabel{prop:monotonicity}
Si $\Gamma \subseteq \Delta$ et $\Gamma \Proves !A$, alors
$\Delta \Proves !A$.
\end{prop}

\begin{proof}
Toute !!{derivation} de~$!A$ à partir de~$\Gamma$ est aussi une
!!{derivation} de~$!A$ à partir de~$\Delta$.
\end{proof}

\begin{prop}[Transitivité]
\ollabel{prop:transitivity}
Si $\Gamma \Proves !A$ et $\{!A\} \cup \Delta \Proves !B$, alors
$\Gamma \cup \Delta \Proves !B$.
\end{prop}

\begin{proof}
Supposons $\{!A\} \cup \Delta \Proves !B$. Il existe alors une
!!{derivation} $!B_1$, \dots, $!B_l = !B$ à partir de
$\{!A\} \cup \Delta$. Certaines étapes de cette !!{derivation} sont
correctes parce qu'une règle renvoie à une ligne antérieure
$!B_i = !A$. Par hypothèse, il existe une !!{derivation} de~$!A$ à
partir de~$\Gamma$, c'est-à-dire une !!{derivation} $!A_1$, \dots,
$!A_k = !A$ dans laquelle chaque $!A_i$ est un axiome, !!a{element}
de~$\Gamma$, ou une étape correcte d'après une règle d'inférence.
Considérons maintenant la suite
\[
!A_1, \dots, !A_k = !A, !B_1, \dots, !B_l = !B.
\]
C'est une !!{derivation} correcte de~$!B$ à partir de
$\Gamma \cup \Delta$ : toute ligne $!B_i = !A$ peut désormais recevoir
la même justification que la ligne~$!A_k = !A$.
\end{proof}

En particulier, si $\Gamma \Proves !A$ et $!A \Proves !B$, alors
$\Gamma \Proves !B$. Il s'ensuit également que, si
$!A_1, \dots, !A_n \Proves !B$ et si $\Gamma \Proves !A_i$ pour
chaque~$i$, alors $\Gamma \Proves !B$.

\begin{prop}
\ollabel{prop:incons}
$\Gamma$ est incohérent si et seulement si
$\Gamma \Proves !A$ pour toute !!{formula}~$!A$.
\end{prop}

\begin{proof}
Exercice.
\end{proof}

\tagprob{FOL}
\begin{prob}
Démontrer la \olref[fol][axd][ptn]{prop:incons}.
\end{prob}
\tagendprob

\tagprob{notFOL}
\begin{prob}
Démontrer la \olref[pl][axd][ptn]{prop:incons}.
\end{prob}
\tagendprob

\begin{prop}[Compacité]
\ollabel{prop:proves-compact}
\begin{enumerate}
\item Si $\Gamma \Proves !A$, alors il existe une partie finie
  $\Gamma_0 \subseteq \Gamma$ telle que $\Gamma_0 \Proves !A$.
\item Si toute partie finie de~$\Gamma$ est cohérente, alors
  $\Gamma$ est cohérent.
\end{enumerate}
\end{prop}

\begin{proof}
\begin{enumerate}
\item Si $\Gamma \Proves !A$, il existe une suite finie de !!{formula}s
  $!A_1$, \dots,~$!A_n$ telle que $!A \ident !A_n$ et que chaque
  $!A_i$ soit un axiome logique, !!a{element} de~$\Gamma$, ou découle
  de !!{formula}s antérieures
  \iftag{FOL}{par une règle d'inférence}{par modus ponens}. Prenons pour
  $\Gamma_0$ l'ensemble des $!A_i$ qui appartiennent à~$\Gamma$.
  La même !!{derivation} est alors une !!{derivation} à partir de
  $\Gamma_0$, et donc $\Gamma_0 \Proves !A$.
\item C'est la contraposée de~(1) dans le cas particulier où
  $!A \ident \lfalse$.
\end{enumerate}
\end{proof}

\iftag{FOL}{\begin{editorial}
Dans la preuve de compacité, la source ne mentionne que les étapes
obtenues par modus ponens. La version avec la balise FOL admet aussi
les règles~\QR. La formulation « par une règle d'inférence » rétablit
la portée commune de la preuve; son argument de finitude est inchangé.
\end{editorial}}{}

\end{document}
